\par\Needspace{6\baselineskip}\medskip
\begin{center}\hypertarget{native-236-article171}{}\pdfbookmark[2]{171.}{article-171}171.\end{center}
To illustrate this investigation, it will be appropriate to take an example from \emph{Vesta}, which has the smallest inclination to the \SourceLeaf{237}ecliptic of all the planets discovered most recently\footnote{Nevertheless, this inclination is still large enough for the orbit to be determined with sufficient confidence and accuracy on the basis of \emph{three} observations: indeed, the first elements, which had been derived in this way from observations only 19 days apart (see \textsc{von Zach} Monatl. Corresp. Vol. XV. p. 595), already come very close to those that will be derived here from four observations 162 days apart.}. We choose the following observations made in Bremen, Paris, Lilienthal and Milan by the distinguished astronomers \textsc{Olbers}, \textsc{Bouvard}, \textsc{Bessel} and \textsc{Oriani}:
\begin{center}
\begingroup\fontsize{9}{11}\selectfont\setlength{\tabcolsep}{3pt}
\begin{tabular}{|l|c|l|}
\hline
\shortstack{Local mean time\\of observation} & Right ascension & Declination\\\hline
1807 March 30. $12^{\mathrm h}33^{\mathrm m}17^{\mathrm s}$ & $183^0\,52'\,40''{,}8$ & $11^0\,54'\,27''$ N.\\
\phantom{1807 }May 17. $8\quad16\quad5$ & $178\quad36\quad42{,}3$ & $11\quad39\quad46{,}8$ —\\
\phantom{1807 }July 11. $10\quad30\quad19$ & $189\quad49\quad7{,}7$ & $3\quad9\quad10{,}1$ —\\
\phantom{1807 }Sept. 8. $7\quad22\quad16$ & $212\quad50\quad3{,}4$ & $8\quad38\quad17{,}0$ S.\\\hline
\end{tabular}\endgroup
\end{center}
\Needspace{12\baselineskip}
For the same times we find from the tables of the Sun's motions
\begin{center}
\begingroup\fontsize{9}{11}\selectfont\setlength{\tabcolsep}{2pt}
\begin{tabular}{|l|c|c|c|c|c|}
\hline
& \shortstack{Sun's long. from\\app. equinox} & Nutation & \shortstack{Distance from\\Earth} & \shortstack{Latitude of\\the Sun} & \shortstack{Apparent obliquity\\of the ecliptic}\\\hline
March 30 & $9^0\,21'\,59''{,}5$ & $+16''{,}8$ & $0{,}999\,6448$ & $+0''{,}23$ & $23^0\,27'\,50''{,}82$\\
May 17 & $55\quad56\quad20{,}0$ & $+16{,}2$ & $1{,}0119789$ & $-0{,}63$ & $49{,}83$\\
July 11 & $108\quad34\quad53{,}3$ & $+17{,}3$ & $1{,}0165795$ & $-0{,}46$ & $49{,}19$\\
Sept. 8 & $165\quad8\quad57{,}1$ & $+16{,}7$ & $1{,}0067421$ & $+0{,}29$ & $49{,}26$\\\hline
\end{tabular}\endgroup
\end{center}

The observed positions of the planet were then converted to longitudes and latitudes using the apparent obliquity of the ecliptic, corrected for nutation and the aberration of the fixed stars, and finally reduced to the beginning of 1807 by removing precession, after which the fictitious positions of the Earth were derived from the positions of the Sun following the instructions of Art. 72 (to account for parallax), and their longitudes were referred to the same epoch by removing nutation and precession; finally, the times were counted from the beginning of the year and reduced to the Paris meridian. This gave the following numbers:
\begin{center}
\begingroup\fontsize{9}{11}\selectfont\setlength{\tabcolsep}{2.5pt}
\begin{tabular}{l|c|c|c|c}
$t,\ t',\ t'',\ t'''$ & $89{,}505\,162$ & $137{,}344502$ & $192{,}419502$ & $251{,}288102$\\
$a,\ a',\ a'',\ a'''$ & $178^0\,43'\,38''{,}87$ & $174^0\,1'\,30''{,}08$ & $187^0\,45'\,42''{,}23$ & $213^0\,34'\,15''{,}63$\\
$\beta,\ \beta',\ \beta'',\ \beta'''$ & $12\quad27\quad6{,}16$ & $10\quad8\quad7{,}80$ & $6\quad47\quad25{,}51$ & $4\quad20\quad21{,}63$\\
$l,\ l',\ l'',\ l'''$ & $189\quad21\quad33{,}71$ & $235\quad56\quad0{,}63$ & $288\quad35\quad20{,}32$ & $345\quad9\quad18{,}69$\\
$\mathrm{logg}\,R,R',R'',R'''$ & $9{,}9997990$ & $0{,}005\,1376$ & $0{,}0071739$ & $0{,}0030625$
\end{tabular}\endgroup
\end{center}

\SourceLeaf{238}From this we deduce
\[
\begin{array}{ll}
\gamma' = 168^{0}32'41''{,}34, & \delta' = 62^{0}23'4''{,}88 \\
\gamma'' = 173\ \ 5\ 15{,}68, & \delta'' = 100\ 45\ 1{,}40 \\
\log a' = 9{,}952\,6104, & b' = -11{,}009\,449 \\
\log a'' = 9{,}999\,4839, & b'' = -\ 2{,}082\,036
\end{array}
\]
\[
\begin{array}{ll}
\varkappa' = -1{,}083\,306, & \log \lambda = 0{,}072\,8800,\\
\log \mu' = 9{,}713\,9702_n & \\
\varkappa'' = +6{,}322\,006, & \log \lambda''' = 0{,}079\,8512_n,\\
\log \mu'' = 9{,}838\,7061 & 
\end{array}
\]
\[
\begin{array}{ll}
A'D = 37^{0}17'51''{,}50, & A''D = 89^{0}24'11''{,}84,\\
\epsilon = 9^{0}5'5''{,}48 & \\
B'D = -25\ \ 5\ 13{,}38, & B''D = -11\ 20\ 49{,}56.
\end{array}
\]
Having completed these preliminary calculations, we proceed to the \emph{first hypothesis}. From the
time intervals we obtain
\[
\begin{array}{l}
\log k(t'-t) = 9{,}915\,3666 \\
\log k(t''-t') = 9{,}976\,5359 \\
\log k(t'''-t'') = 0{,}005\,4651
\end{array}
\]
and from these the first approximate values
\[
\begin{array}{ll}
\log P' = 0{,}06117, & \log(1+P') = 0{,}33269,\\
\log Q' = 9{,}59087 & \\
\log P'' = 9{,}97107, & \log(1+P'') = 0{,}28681,\\
\log Q'' = 9{,}68097, & 
\end{array}
\]
and further from these
\[
\begin{array}{ll}
c' = -7{,}68361, & \log d' = 0{,}04666_n \\
c'' = +2{,}20771, & \log d'' = 0{,}12552.
\end{array}
\]
With these values, after a few trials, the following solution of equations I. II
is obtained:
\[
\begin{array}{ll}
x' = 2{,}04856, & z' = 23^{0}38'17'',\\
\log r' = 0{,}34951 & \\
x'' = 1{,}95745, & z'' = 27\ \ 2\ \ 0,\\
\log r'' = 0{,}34194. & 
\end{array}
\]
From $z'$, $z''$ and $\epsilon$ we obtain $C'C'' = v''-v' = 17^{0}7'5''$: from this
$v'-v$, $r$, $v'''-v''$, $r'''$ are to be determined by the following equations:
\[
\begin{array}{l}
\log r\ \sin(v'-v) = 9{,}74962, \\ \log r\ \sin(v'-v+17^{0}7'5'') = 0{,}07500 \\
\log r'''\sin(v'''-v'') = 9{,}84729, \\ \log r'''\sin(v'''-v'+17^{0}7'5'') = 0{,}10733,
\end{array}
\]
from which we obtain
\SourceLeaf{239}
\[
\begin{array}{ll}
v'-v = 14^{0}14'32'', & \log r = 0{,}35865 \\
v'''-v'' = 18\ 48\ 33, & \log r''' = 0{,}33887.
\end{array}
\]
Finally we find
\[
\begin{gathered}
\log(\eta 01) = 0{,}00426, \\ \log(\eta 12) = 0{,}00599, \\ \log(\eta 23) = 0{,}00711,
\end{gathered}
\]
and from this the corrected values of $P'$, $P''$, $Q'$, $Q''$:
\[
\begin{array}{ll}
\log P' = 0{,}05944, & \log Q' = 9{,}60374 \\
\log P'' = 9{,}97219, & \log Q'' = 9{,}69581,
\end{array}
\]
on which the \emph{second hypothesis} is to be based. Its principal results
are as follows:
\[
\begin{array}{ll}
c' = -7{,}67820, & \log d' = 0{,}045\,736_n\mathord{\rule[.1ex]{.35em}{.35pt}} \\
c'' = +2{,}21061, & \log d'' = 0{,}126\,054
\end{array}
\]
\[
\begin{array}{ll}
x' = 2{,}03308, & z' = 23^{0}47'54'',\\
\log r' = 0{,}346\,747 & \\
x'' = 1{,}94290, & z'' = 27\ 12\ 25,\\
\log r'' = 0{,}339\,373 & 
\end{array}
\]
\[
C'C'' = v''-v' = 17^{0}8'0''
\]
\[
\begin{array}{ll}
v'-v = 14^{0}21'36'', & \log r = 0{,}354\,687 \\
v'''-v'' = 18\ 50\ 43, & \log r''' = 0{,}334\,564
\end{array}
\]
\[
\begin{gathered}
\log(\eta 01) = 0{,}004\,359, \\ \log(\eta 12) = 0{,}006\,102, \\ \log(\eta 23) = 0{,}007\,280.
\end{gathered}
\]
This gives the values of $P'$, $P''$, $Q'$, $Q''$ corrected once more:
\[
\begin{array}{ll}
\log P' = 0{,}059\,426, & \log Q' = 9{,}604\,749 \\
\log P'' = 9{,}972\,249, & \log Q'' = 9{,}697\,564,
\end{array}
\]
if we use these to proceed to the \emph{third hypothesis}, the following numbers
result:
\[
\begin{array}{ll}
c' = -7{,}67815, & \log d' = 0{,}045\,729_n \\
c'' = +2{,}21076, & \log d'' = 0{,}126\,082
\end{array}
\]
\[
\begin{array}{ll}
x' = 2{,}03255, & z' = 23^{0}48'14'',\\
\log r' = 0{,}346\,653 & \\
x'' = 1{,}94235, & z'' = 27\ 12\ 49,\\
\log r'' = 0{,}339\,276 & 
\end{array}
\]
\[
C'C'' = v''-v' = 17^{0}8'4''
\]
\[
\begin{array}{ll}
v'-v = 14^{0}21'49'', & \log r = 0{,}354\,522 \\
v'''-v'' = 18\ 51\ \ 7, & \log r''' = 0{,}334\,290
\end{array}
\]
\[
\begin{gathered}
\log(\eta 01) = 0{,}004\,363, \\ \log(\eta 12) = 0{,}006\,106, \\ \log(\eta 23) = 0{,}007\,290.
\end{gathered}
\]
\SourceLeaf{240}
If the distances from the Earth are now calculated following the instructions of the preceding article,
we obtain:
\[
\begin{array}{ll}
\rho' = 1{,}5635, & \rho'' = 2{,}1319 \\
\log \rho \cos\beta = 0{,}09876, & \log \rho'''\cos\beta''' = 0{,}42842 \\
\log \rho \sin\beta = 9{,}44252, & \log \rho'''\sin\beta''' = 9{,}30905 \\
\beta = 12^{0}26'40'', & \beta''' = 4^{0}20'39'' \\
\log \rho = 0{,}10909, & \log \rho''' = 0{,}42967.
\end{array}
\]
From this we find
\[
\begin{array}{c|c|c}
 & \text{Corrections to the times} & \text{Corrected times} \\
\hline
\text{I} & 0{,}007\,335 & 89{,}497\,827 \\
\text{II} & 0{,}008\,921 & 137{,}335\,581 \\
\text{III} & 0{,}012\,165 & 192{,}407\,337 \\
\text{IV} & 0{,}015\,346 & 251{,}272\,756
\end{array}
\]
which gives the values of the quantities $P'$, $P''$, $Q'$, $Q''$ corrected once more
\[
\begin{array}{ll}
\log P' = 0{,}059\,415, & \log Q' = 9{,}604\,782 \\
\log P'' = 9{,}972\,253, & \log Q'' = 9{,}697\,687.
\end{array}
\]
Finally, if the \emph{fourth hypothesis} is formed with these new values, the following numbers
result:
\[
\begin{array}{ll}
c' = -7{,}678\,116, & \log d' = 0{,}045\,723_n \\
c'' = +2{,}210\,773, & \log d'' = 0{,}126\,084
\end{array}
\]
\[
\begin{array}{ll}
x' = 2{,}032\,473, & z' = 23^{0}48'16''{,}7,\\
\log r' = 0{,}346\,638 & \\
x'' = 1{,}942\,281, & z'' = 27\ 12\ 51{,}7,\\
\log r'' = 0{,}339\,263 & 
\end{array}
\]
\[
\begin{gathered}
v''-v' = 17^{0}8'5''{,}1, \\ \tfrac{1}{2}(u''+u') = 176^{0}7'50''{,}5, \\
\tfrac{1}{2}(u''-u') = 4^{0}33'23''{,}6
\end{gathered}
\]
\[
\begin{array}{ll}
v'-v = 14^{0}21'51''{,}9, & \log r = 0{,}354\,503 \\
v'''-v'' = 18\ 51\ \ 9{,}5, & \log r''' = 0{,}334\,263.
\end{array}
\]
These numbers differ so little from those supplied by the third hypothesis
that we can now safely proceed to determine the elements themselves.
We first determine the orientation of the orbital plane. Following the instructions of Art.\ 149
we find, from $\gamma'$, $u'$ and $A'C' = \delta'-z'$, the inclination of the orbit
$= 7^{0}8'14''{,}8$, the longitude of the ascending node $103^{0}16'37''{,}2$, and the argument
of latitude in the second observation $94^{0}36'4''{,}9$, \SourceLeaf{241}and hence the longitude in the orbit $197^{0}52'42''{,}1$; likewise from $\gamma''$, $u''$
and $A''C'' = \delta''-z''$ we obtain the inclination of the orbit $= 7^{0}8'14''{,}8$,
the longitude of the ascending node $103^{0}16'37''{,}5$, and the argument of latitude in
the third observation $111^{0}44'9''{,}7$, and hence the longitude in the orbit
$215^{0}0'47''{,}2$. Hence the longitude in the orbit for the first observation will be
$183^{0}30'50''{,}2$, and for the fourth $233^{0}51'56''{,}7$. If the dimensions of the orbit are now determined from $t'''-t$,
$r$, $r'''$ and $v'''-v = 50^{0}21'6''{,}5$,
we obtain:
\begin{center}
\begingroup\fontsize{9}{11}\selectfont\setlength{\tabcolsep}{3pt}
\begin{tabular}{@{}lr@{}}
True anomaly for the first position . . . . . . . . . . . . . . . . . & $293^{0}33'43''{,}7$ \\
True anomaly for the fourth position . . . . . . . . . . . . . . . . . & $343\ \ 54\ \ 50{,}2$ \\
Hence the longitude of perihelion . . . . . . . . . . . . . . . . . . . . . . . & $249\ \ 57\ \ \ 6{,}5$ \\
Mean anomaly for the first position . . . . . . . . . . . . . . . . & $302\ \ 33\ \ 32{,}6$ \\
Mean anomaly for the fourth position . . . . . . . . . . . . . . . & $346\ \ 32\ \ 25{,}2$ \\
Mean daily sidereal motion . . . . . . . . . . . . . . . . . & $978''{,}7216$ \\
Mean anomaly for the beginning of 1807 . . . . . . . . . . . & $278^{0}13'39''{,}1$ \\
Mean longitude for the same epoch . . . . . . . . . . . . & $168\ \ 10\ \ 45{,}6$ \\
Angle $\varphi$ . . . . . . . . . . . . . . . . . . . . . . . . . . . . . . . . . . & $5\ \ \ 2\ \ 58{,}1$ \\
Logarithm of the semimajor axis . . . . . . . . . . . . . . . . . . & $0{,}372\,898.$
\end{tabular}\endgroup
\end{center}

If the geocentric positions of the planet are calculated from these elements for the corrected times $t$, $t'$, $t''$, $t'''$,
the four longitudes agree with $a$, $a'$,
$a''$, $a'''$, and the two intermediate latitudes with $\beta'$, $\beta''$, to within
one tenth of a second of arc; the first and last latitudes, however,
come out as $12^{0}26'43''{,}7$ and $4^{0}20'40''{,}1$, the former too low by $22''{,}4$
and the latter too high by $18''{,}5$. However, if the remaining elements are held fixed
and only the inclination of the orbit is increased by $6''$ and the longitude of the node decreased by $4' 40''$,
the errors, distributed among all the latitudes, will be reduced to a few
seconds of arc, and the longitudes will be affected by only very slight errors,
which will themselves be reduced almost to nothing if, in addition,
the epoch of longitude is decreased by $2''$.
\par\medskip\begin{center}\rule{.24\linewidth}{.35pt}\end{center}
